% --- Archon physics formalization source begin ---
% archon:physics
% archon:covers PhyXMiniProblems/problem_phyx_mini_0292.lean
% archon:source-report reports/phyx_mini/problem_phyx_mini_0292.source.json

\chapter{Physics problem phyx\_mini\_0292}
\label{ch:PhyXMiniProblems_problem_phyx_mini_0292}

\paragraph{Problem source.}
\#\# Physical scenario

In figure, string 1 has a linear density of \textbackslash{}(3.00\textbackslash{},\textbackslash{}mathrm\{g/m\}\textbackslash{}), and string 2 has a linear density of \textbackslash{}(5.00\textbackslash{},\textbackslash{}mathrm\{g/m\}\textbackslash{}). They are under tension due to the hanging block of mass \textbackslash{}(M = 500\textbackslash{},\textbackslash{}mathrm\{g\}\textbackslash{}).

\#\# Question

Calculate the wave speed on string 2.

\#\# Answer choices

- A: A: 19.8 m/s

- B: B: 20.3 m/s

- C: C: 21.4 m/s

- D: D: 22.1 m/s

\#\# Figure caption (auxiliary; use the image as primary evidence)

The image depicts a pulley system involving two strings labeled as "String 1" and "String 2." The strings pass over two pulleys anchored to a fixed structure. There is a knot joining String 1 and String 2. Hanging from the knot is a single string supporting a mass labeled as "M." The setup shows how the strings and pulleys are arranged to handle the mass, with the pulleys facilitating movement.

\#\# Dataset metadata

Wave Properties; Physical Model Grounding Reasoning, Multi-Formula Reasoning

\paragraph{Recorded answer/context.}
D: D: 22.1 m/s

\paragraph{Figure/image path.}
/root/proposal\_for\_physic/science-mango/phyx\_mini\_run/phyx\_data/test\_image/292.png

\paragraph{Formalization target.}
create a compiling Lean file with sorry bodies at `PhyXMiniProblems/problem\_phyx\_mini\_0292.lean`. The Lean declarations must preserve the physical quantities, dimensions or dimensional roles, figure labels, governing-law hypotheses, and final relation expressed by this problem.
Use Mathlib/Physlib names found through LeanExplore where available. If a physics API is missing, introduce faithful local abstractions rather than scalar placeholder aliases.

\begin{theorem}[Physics formalization target]
\label{thm:physics:phyx_mini_0292:target}
\lean{PhyXMiniProblems.ProblemPhyXMini0292.problem_phyx_mini_0292}
\uses{def:physics:phyx-mini-0292:phyxminiproblems-problemphyxmini0292-pulleystringsetup, def:physics:phyx-mini-0292:phyxminiproblems-problemphyxmini0292-matchesproblemandprimaryfigure, def:physics:phyx-mini-0292:phyxminiproblems-problemphyxmini0292-usesstandardearthgravity, def:physics:phyx-mini-0292:phyxminiproblems-problemphyxmini0292-haspositivephysicalparameters, def:physics:phyx-mini-0292:phyxminiproblems-problemphyxmini0292-satisfiesidealpulleyknotandwavelaws, def:physics:phyx-mini-0292:phyxminiproblems-problemphyxmini0292-recordedanswerchoice, def:physics:phyx-mini-0292:phyxminiproblems-problemphyxmini0292-matchesanswerchoice}
The assigned autoformalize agent should translate this physics problem into Lean declarations in the covered file, with theorem and lemma proof bodies written as `by sorry`.
\end{theorem}
\begin{proof}
This is an autoformalization task, not a proof task. Produce statements that can later be proved by the physics prover without weakening the physical model.
\end{proof}

% --- Archon declaration topology begin ---
\section{Declaration topology}
\begin{definition}[Mass Quantity]
  \label{def:physics:phyx-mini-0292:phyxminiproblems-problemphyxmini0292-massquantity}
  \lean{PhyXMiniProblems.ProblemPhyXMini0292.MassQuantity}
  A nonnegative physical mass
\end{definition}
\begin{proof}
  This is the declared model definition of Mass Quantity.
\end{proof}
\begin{definition}[Acceleration Quantity]
  \label{def:physics:phyx-mini-0292:phyxminiproblems-problemphyxmini0292-accelerationquantity}
  \lean{PhyXMiniProblems.ProblemPhyXMini0292.AccelerationQuantity}
  A nonnegative physical acceleration magnitude
\end{definition}
\begin{proof}
  This is the declared model definition of Acceleration Quantity.
\end{proof}
\begin{definition}[Tension Quantity]
  \label{def:physics:phyx-mini-0292:phyxminiproblems-problemphyxmini0292-tensionquantity}
  \lean{PhyXMiniProblems.ProblemPhyXMini0292.TensionQuantity}
  A nonnegative force magnitude, used here for a string-segment tension
\end{definition}
\begin{proof}
  This is the declared model definition of Tension Quantity.
\end{proof}
\begin{definition}[Linear Mass Density Quantity]
  \label{def:physics:phyx-mini-0292:phyxminiproblems-problemphyxmini0292-linearmassdensityquantity}
  \lean{PhyXMiniProblems.ProblemPhyXMini0292.LinearMassDensityQuantity}
  A nonnegative mass per unit length of string
\end{definition}
\begin{proof}
  This is the declared model definition of Linear Mass Density Quantity.
\end{proof}
\begin{definition}[Speed Quantity]
  \label{def:physics:phyx-mini-0292:phyxminiproblems-problemphyxmini0292-speedquantity}
  \lean{PhyXMiniProblems.ProblemPhyXMini0292.SpeedQuantity}
  A nonnegative propagation speed
\end{definition}
\begin{proof}
  This is the declared model definition of Speed Quantity.
\end{proof}
\begin{definition}[quantity Readout]
  \label{def:physics:phyx-mini-0292:phyxminiproblems-problemphyxmini0292-quantityreadout}
  \lean{PhyXMiniProblems.ProblemPhyXMini0292.quantityReadout}
  Read a nonnegative dimensionful quantity in a coherent unit system
\end{definition}
\begin{proof}
  This is the declared model definition of quantity Readout.
\end{proof}
\begin{definition}[mass In Kilograms]
  \label{def:physics:phyx-mini-0292:phyxminiproblems-problemphyxmini0292-massinkilograms}
  \lean{PhyXMiniProblems.ProblemPhyXMini0292.massInKilograms}
  \uses{def:physics:phyx-mini-0292:phyxminiproblems-problemphyxmini0292-massquantity, def:physics:phyx-mini-0292:phyxminiproblems-problemphyxmini0292-quantityreadout}
  Kilogram readout of a mass
\end{definition}
\begin{proof}
  This is the declared model definition of mass In Kilograms.
\end{proof}
\begin{definition}[acceleration In Meters Per Second Squared]
  \label{def:physics:phyx-mini-0292:phyxminiproblems-problemphyxmini0292-accelerationinmeterspersecondsquared}
  \lean{PhyXMiniProblems.ProblemPhyXMini0292.accelerationInMetersPerSecondSquared}
  \uses{def:physics:phyx-mini-0292:phyxminiproblems-problemphyxmini0292-accelerationquantity, def:physics:phyx-mini-0292:phyxminiproblems-problemphyxmini0292-quantityreadout}
  Meter-per-second-squared readout of an acceleration
\end{definition}
\begin{proof}
  This is the declared model definition of acceleration In Meters Per Second Squared.
\end{proof}
\begin{definition}[tension In Newtons]
  \label{def:physics:phyx-mini-0292:phyxminiproblems-problemphyxmini0292-tensioninnewtons}
  \lean{PhyXMiniProblems.ProblemPhyXMini0292.tensionInNewtons}
  \uses{def:physics:phyx-mini-0292:phyxminiproblems-problemphyxmini0292-tensionquantity, def:physics:phyx-mini-0292:phyxminiproblems-problemphyxmini0292-quantityreadout}
  Newton readout of a tension
\end{definition}
\begin{proof}
  This is the declared model definition of tension In Newtons.
\end{proof}
\begin{definition}[linear Mass Density In Kilograms Per Meter]
  \label{def:physics:phyx-mini-0292:phyxminiproblems-problemphyxmini0292-linearmassdensityinkilogramspermeter}
  \lean{PhyXMiniProblems.ProblemPhyXMini0292.linearMassDensityInKilogramsPerMeter}
  \uses{def:physics:phyx-mini-0292:phyxminiproblems-problemphyxmini0292-linearmassdensityquantity, def:physics:phyx-mini-0292:phyxminiproblems-problemphyxmini0292-quantityreadout}
  Kilogram-per-meter readout of a linear mass density
\end{definition}
\begin{proof}
  This is the declared model definition of linear Mass Density In Kilograms Per Meter.
\end{proof}
\begin{definition}[speed In Meters Per Second]
  \label{def:physics:phyx-mini-0292:phyxminiproblems-problemphyxmini0292-speedinmeterspersecond}
  \lean{PhyXMiniProblems.ProblemPhyXMini0292.speedInMetersPerSecond}
  \uses{def:physics:phyx-mini-0292:phyxminiproblems-problemphyxmini0292-speedquantity, def:physics:phyx-mini-0292:phyxminiproblems-problemphyxmini0292-quantityreadout}
  Meter-per-second readout of a wave speed
\end{definition}
\begin{proof}
  This is the declared model definition of speed In Meters Per Second.
\end{proof}
\begin{definition}[String Label]
  \label{def:physics:phyx-mini-0292:phyxminiproblems-problemphyxmini0292-stringlabel}
  \lean{PhyXMiniProblems.ProblemPhyXMini0292.StringLabel}
  The two strings explicitly labelled in the primary figure
\end{definition}
\begin{proof}
  This is the declared model definition of String Label.
\end{proof}
\begin{definition}[Pulley Label]
  \label{def:physics:phyx-mini-0292:phyxminiproblems-problemphyxmini0292-pulleylabel}
  \lean{PhyXMiniProblems.ProblemPhyXMini0292.PulleyLabel}
  The three separately drawn pulleys
\end{definition}
\begin{proof}
  This is the declared model definition of Pulley Label.
\end{proof}
\begin{definition}[Pulley Mounting]
  \label{def:physics:phyx-mini-0292:phyxminiproblems-problemphyxmini0292-pulleymounting}
  \lean{PhyXMiniProblems.ProblemPhyXMini0292.PulleyMounting}
  Whether a pulley is anchored to the support or moves with the load
\end{definition}
\begin{proof}
  This is the declared model definition of Pulley Mounting.
\end{proof}
\begin{definition}[Figure Component]
  \label{def:physics:phyx-mini-0292:phyxminiproblems-problemphyxmini0292-figurecomponent}
  \lean{PhyXMiniProblems.ProblemPhyXMini0292.FigureComponent}
  Individually identifiable objects shown in the primary figure
\end{definition}
\begin{proof}
  This is the declared model definition of Figure Component.
\end{proof}
\begin{definition}[Endpoint Attachment]
  \label{def:physics:phyx-mini-0292:phyxminiproblems-problemphyxmini0292-endpointattachment}
  \lean{PhyXMiniProblems.ProblemPhyXMini0292.EndpointAttachment}
  Objects to which the support-side and knot-side string ends attach
\end{definition}
\begin{proof}
  This is the declared model definition of Endpoint Attachment.
\end{proof}
\begin{definition}[String Endpoint]
  \label{def:physics:phyx-mini-0292:phyxminiproblems-problemphyxmini0292-stringendpoint}
  \lean{PhyXMiniProblems.ProblemPhyXMini0292.StringEndpoint}
  The two distinguished ends of either pictured string
\end{definition}
\begin{proof}
  This is the declared model definition of String Endpoint.
\end{proof}
\begin{definition}[String Route]
  \label{def:physics:phyx-mini-0292:phyxminiproblems-problemphyxmini0292-stringroute}
  \lean{PhyXMiniProblems.ProblemPhyXMini0292.StringRoute}
  The full route of each string through the pictured pulley system
\end{definition}
\begin{proof}
  This is the declared model definition of String Route.
\end{proof}
\begin{definition}[String Segment]
  \label{def:physics:phyx-mini-0292:phyxminiproblems-problemphyxmini0292-stringsegment}
  \lean{PhyXMiniProblems.ProblemPhyXMini0292.StringSegment}
  The five straight string portions separated by pulleys and by the knot. The two String-1 portions adjacent to the movable pulley supply its two upward supporting tensions
\end{definition}
\begin{proof}
  This is the declared model definition of String Segment.
\end{proof}
\begin{definition}[string Label]
  \label{def:physics:phyx-mini-0292:phyxminiproblems-problemphyxmini0292-stringsegment-stringlabel}
  \lean{PhyXMiniProblems.ProblemPhyXMini0292.StringSegment.stringLabel}
  \uses{def:physics:phyx-mini-0292:phyxminiproblems-problemphyxmini0292-stringlabel, def:physics:phyx-mini-0292:phyxminiproblems-problemphyxmini0292-stringsegment}
  The string to which a named straight segment belongs
\end{definition}
\begin{proof}
  This is the declared model definition of string Label.
\end{proof}
\begin{definition}[Pulley String Setup]
  \label{def:physics:phyx-mini-0292:phyxminiproblems-problemphyxmini0292-pulleystringsetup}
  \lean{PhyXMiniProblems.ProblemPhyXMini0292.PulleyStringSetup}
  \uses{def:physics:phyx-mini-0292:phyxminiproblems-problemphyxmini0292-massquantity, def:physics:phyx-mini-0292:phyxminiproblems-problemphyxmini0292-accelerationquantity, def:physics:phyx-mini-0292:phyxminiproblems-problemphyxmini0292-tensionquantity, def:physics:phyx-mini-0292:phyxminiproblems-problemphyxmini0292-linearmassdensityquantity, def:physics:phyx-mini-0292:phyxminiproblems-problemphyxmini0292-speedquantity, def:physics:phyx-mini-0292:phyxminiproblems-problemphyxmini0292-stringlabel, def:physics:phyx-mini-0292:phyxminiproblems-problemphyxmini0292-pulleylabel, def:physics:phyx-mini-0292:phyxminiproblems-problemphyxmini0292-pulleymounting, def:physics:phyx-mini-0292:phyxminiproblems-problemphyxmini0292-figurecomponent, def:physics:phyx-mini-0292:phyxminiproblems-problemphyxmini0292-endpointattachment, def:physics:phyx-mini-0292:phyxminiproblems-problemphyxmini0292-stringendpoint, def:physics:phyx-mini-0292:phyxminiproblems-problemphyxmini0292-stringroute, def:physics:phyx-mini-0292:phyxminiproblems-problemphyxmini0292-stringsegment}
  Independent physical quantities and qualitative figure data. In particular, neither the tension on String 2 nor its requested wave speed is assigned a numerical value here
\end{definition}
\begin{proof}
  This is the declared model definition of Pulley String Setup.
\end{proof}
\begin{definition}[Matches Problem And Primary Figure]
  \label{def:physics:phyx-mini-0292:phyxminiproblems-problemphyxmini0292-matchesproblemandprimaryfigure}
  \lean{PhyXMiniProblems.ProblemPhyXMini0292.MatchesProblemAndPrimaryFigure}
  \uses{def:physics:phyx-mini-0292:phyxminiproblems-problemphyxmini0292-massinkilograms, def:physics:phyx-mini-0292:phyxminiproblems-problemphyxmini0292-linearmassdensityinkilogramspermeter, def:physics:phyx-mini-0292:phyxminiproblems-problemphyxmini0292-pulleystringsetup}
  Figure topology and numerical readouts supplied by the problem. The image, taken as primary evidence, places M below the lower movable pulley rather than directly below the knot
\end{definition}
\begin{proof}
  This is the declared model definition of Matches Problem And Primary Figure.
\end{proof}
\begin{definition}[Uses Standard Earth Gravity]
  \label{def:physics:phyx-mini-0292:phyxminiproblems-problemphyxmini0292-usesstandardearthgravity}
  \lean{PhyXMiniProblems.ProblemPhyXMini0292.UsesStandardEarthGravity}
  \uses{def:physics:phyx-mini-0292:phyxminiproblems-problemphyxmini0292-accelerationinmeterspersecondsquared, def:physics:phyx-mini-0292:phyxminiproblems-problemphyxmini0292-pulleystringsetup}
  The conventional near-Earth gravitational readout used by the choices
\end{definition}
\begin{proof}
  This is the declared model definition of Uses Standard Earth Gravity.
\end{proof}
\begin{definition}[Has Positive Physical Parameters]
  \label{def:physics:phyx-mini-0292:phyxminiproblems-problemphyxmini0292-haspositivephysicalparameters}
  \lean{PhyXMiniProblems.ProblemPhyXMini0292.HasPositivePhysicalParameters}
  \uses{def:physics:phyx-mini-0292:phyxminiproblems-problemphyxmini0292-massinkilograms, def:physics:phyx-mini-0292:phyxminiproblems-problemphyxmini0292-accelerationinmeterspersecondsquared, def:physics:phyx-mini-0292:phyxminiproblems-problemphyxmini0292-tensioninnewtons, def:physics:phyx-mini-0292:phyxminiproblems-problemphyxmini0292-linearmassdensityinkilogramspermeter, def:physics:phyx-mini-0292:phyxminiproblems-problemphyxmini0292-speedinmeterspersecond, def:physics:phyx-mini-0292:phyxminiproblems-problemphyxmini0292-pulleystringsetup}
  Positivity and nondegeneracy conditions for the taut-string model
\end{definition}
\begin{proof}
  This is the declared model definition of Has Positive Physical Parameters.
\end{proof}
\begin{definition}[Satisfies Ideal Pulley Knot And Wave Laws]
  \label{def:physics:phyx-mini-0292:phyxminiproblems-problemphyxmini0292-satisfiesidealpulleyknotandwavelaws}
  \lean{PhyXMiniProblems.ProblemPhyXMini0292.SatisfiesIdealPulleyKnotAndWaveLaws}
  \uses{def:physics:phyx-mini-0292:phyxminiproblems-problemphyxmini0292-massinkilograms, def:physics:phyx-mini-0292:phyxminiproblems-problemphyxmini0292-accelerationinmeterspersecondsquared, def:physics:phyx-mini-0292:phyxminiproblems-problemphyxmini0292-tensioninnewtons, def:physics:phyx-mini-0292:phyxminiproblems-problemphyxmini0292-linearmassdensityinkilogramspermeter, def:physics:phyx-mini-0292:phyxminiproblems-problemphyxmini0292-speedinmeterspersecond, def:physics:phyx-mini-0292:phyxminiproblems-problemphyxmini0292-pulleystringsetup}
  The textbook idealizations needed for the calculation: negligible string weight in the static balance and frictionless ideal pulleys transmit one tension along every segment of a given string; the two upward String-1 segments balance the lower pulley and mass; the massless static knot balances the adjacent String-1 and String-2 pulls; every taut string obeys v = sqrt (T / mu). These laws are stated generically in terms of independent setup quantities. They contain neither sqrt 490 nor the displayed answer 22.1 m/s
\end{definition}
\begin{proof}
  This is the declared model definition of Satisfies Ideal Pulley Knot And Wave Laws.
\end{proof}
\begin{lemma}[string One Supporting Tensions are 49 over 20 newtons]
  \label{lem:physics:phyx-mini-0292:phyxminiproblems-problemphyxmini0292-stringonesupportingtensions-are-49-over-20-newtons}
  \lean{PhyXMiniProblems.ProblemPhyXMini0292.stringOneSupportingTensions_are_49_over_20_newtons}
  \uses{def:physics:phyx-mini-0292:phyxminiproblems-problemphyxmini0292-tensioninnewtons, def:physics:phyx-mini-0292:phyxminiproblems-problemphyxmini0292-pulleystringsetup, def:physics:phyx-mini-0292:phyxminiproblems-problemphyxmini0292-matchesproblemandprimaryfigure, def:physics:phyx-mini-0292:phyxminiproblems-problemphyxmini0292-usesstandardearthgravity, def:physics:phyx-mini-0292:phyxminiproblems-problemphyxmini0292-satisfiesidealpulleyknotandwavelaws}
  Each String-1 segment supporting the movable pulley has tension 2.45 N
\end{lemma}
\begin{proof}
  The conclusion follows from the governing relations and figure readouts named in the statement of string One Supporting Tensions are 49 over 20 newtons.
\end{proof}
\begin{lemma}[string Two Tension is 49 over 20 newtons]
  \label{lem:physics:phyx-mini-0292:phyxminiproblems-problemphyxmini0292-stringtwotension-is-49-over-20-newtons}
  \lean{PhyXMiniProblems.ProblemPhyXMini0292.stringTwoTension_is_49_over_20_newtons}
  \uses{def:physics:phyx-mini-0292:phyxminiproblems-problemphyxmini0292-tensioninnewtons, def:physics:phyx-mini-0292:phyxminiproblems-problemphyxmini0292-pulleystringsetup, def:physics:phyx-mini-0292:phyxminiproblems-problemphyxmini0292-matchesproblemandprimaryfigure, def:physics:phyx-mini-0292:phyxminiproblems-problemphyxmini0292-usesstandardearthgravity, def:physics:phyx-mini-0292:phyxminiproblems-problemphyxmini0292-satisfiesidealpulleyknotandwavelaws}
  Static balance at the massless knot gives String 2 the same 2.45 N tension
\end{lemma}
\begin{proof}
  The conclusion follows from the governing relations and figure readouts named in the statement of string Two Tension is 49 over 20 newtons.
\end{proof}
\begin{lemma}[string Two Wave Speed exact]
  \label{lem:physics:phyx-mini-0292:phyxminiproblems-problemphyxmini0292-stringtwowavespeed-exact}
  \lean{PhyXMiniProblems.ProblemPhyXMini0292.stringTwoWaveSpeed_exact}
  \uses{def:physics:phyx-mini-0292:phyxminiproblems-problemphyxmini0292-speedinmeterspersecond, def:physics:phyx-mini-0292:phyxminiproblems-problemphyxmini0292-pulleystringsetup, def:physics:phyx-mini-0292:phyxminiproblems-problemphyxmini0292-matchesproblemandprimaryfigure, def:physics:phyx-mini-0292:phyxminiproblems-problemphyxmini0292-usesstandardearthgravity, def:physics:phyx-mini-0292:phyxminiproblems-problemphyxmini0292-haspositivephysicalparameters, def:physics:phyx-mini-0292:phyxminiproblems-problemphyxmini0292-satisfiesidealpulleyknotandwavelaws}
  With mu\_2 = 0.005 kg/m, the unrounded requested speed is sqrt ((2.45 N) / (0.005 kg/m)) = sqrt 490 m/s
\end{lemma}
\begin{proof}
  The conclusion follows from the governing relations and figure readouts named in the statement of string Two Wave Speed exact.
\end{proof}
\begin{definition}[Answer Choice]
  \label{def:physics:phyx-mini-0292:phyxminiproblems-problemphyxmini0292-answerchoice}
  \lean{PhyXMiniProblems.ProblemPhyXMini0292.AnswerChoice}
  Labels printed beside the four multiple-choice answers
\end{definition}
\begin{proof}
  This is the declared model definition of Answer Choice.
\end{proof}
\begin{definition}[speed In Meters Per Second]
  \label{def:physics:phyx-mini-0292:phyxminiproblems-problemphyxmini0292-answerchoice-speedinmeterspersecond}
  \lean{PhyXMiniProblems.ProblemPhyXMini0292.AnswerChoice.speedInMetersPerSecond}
  \uses{def:physics:phyx-mini-0292:phyxminiproblems-problemphyxmini0292-speedinmeterspersecond, def:physics:phyx-mini-0292:phyxminiproblems-problemphyxmini0292-answerchoice}
  Displayed wave-speed readout for each choice, in meters per second
\end{definition}
\begin{proof}
  This is the declared model definition of speed In Meters Per Second.
\end{proof}
\begin{definition}[recorded Answer Choice]
  \label{def:physics:phyx-mini-0292:phyxminiproblems-problemphyxmini0292-recordedanswerchoice}
  \lean{PhyXMiniProblems.ProblemPhyXMini0292.recordedAnswerChoice}
  \uses{def:physics:phyx-mini-0292:phyxminiproblems-problemphyxmini0292-answerchoice}
  The answer label recorded in the supplied dataset metadata
\end{definition}
\begin{proof}
  This is the declared model definition of recorded Answer Choice.
\end{proof}
\begin{definition}[Matches Answer Choice]
  \label{def:physics:phyx-mini-0292:phyxminiproblems-problemphyxmini0292-matchesanswerchoice}
  \lean{PhyXMiniProblems.ProblemPhyXMini0292.MatchesAnswerChoice}
  \uses{def:physics:phyx-mini-0292:phyxminiproblems-problemphyxmini0292-speedinmeterspersecond, def:physics:phyx-mini-0292:phyxminiproblems-problemphyxmini0292-pulleystringsetup, def:physics:phyx-mini-0292:phyxminiproblems-problemphyxmini0292-answerchoice}
  Agreement with a speed displayed to one decimal place. The tolerance is 0.05 m/s, half a unit in the last displayed digit
\end{definition}
\begin{proof}
  This is the declared model definition of Matches Answer Choice.
\end{proof}
% --- Archon declaration topology end ---
% --- Archon physics formalization source end ---
